RBM reframes what is often called ``sound healing'' as a more specific claim about biophysical entrainment: acoustic and vibrotactile inputs are not treated as vague wellness signals, but as structured perturbations capable of coupling to measurable physiological rhythms. In this view, the relevant question is not whether sound is ``healing'' in a general sense, but which frequencies, amplitudes, phase relationships, and delivery modes can reliably bias neural, autonomic, or metabolic dynamics toward a desired state.

This framing also clarifies the role of symbolic medicine. Subjective meaning, expectation, ritual, and imagery are not dismissed as epiphenomenal; rather, they are treated as part of the interface through which a stimulus is interpreted and routed into the body. Symbolic forms can therefore bridge lived experience and real modulation by shaping attention, autonomic tone, and compliance with a stimulus protocol. RBM does not require a separation between ``mind'' and ``mechanism'' so much as a disciplined account of how meaning can participate in physiological change.

At a broader level, RBM suggests a progression from perceptual grammar tuning to reality interface engineering. If perception is partly organized by learned patterns of rhythm, salience, and expectation, then targeted sonic and vibrotactile design may alter the grammar by which experience is parsed and stabilized. The practical implication is that frequency-based interventions may be used not only to shift isolated biomarkers, but to reconfigure the coupling between sensation, interpretation, and embodied response. In that sense, RBM positions sound as a tool for perceptual engineering: a way to shape the interface through which organisms encounter and regulate their world.
